% [原文 p.57]
$K$ 与 $DK$ 的互补维数以外的闭链的相交，也显然有几何意义。于是「闭链的相交」（以及由此得到的同调类的相交）这一运算，确实可以看作来自真正的几何相交。

\section{相对同调。对的正合列。同调公理与 CW 复形}
\takeaway{对复形对 $(K,L)$ 取商链复形 $C(K)/C(L)$，即得相对同调与相对上同调；它们通过长正合列 (3.3)、(3.4) 与 $K$、$L$ 的（上）同调相连。同调公理把这些性质公理化，CW 复形则把计算归约为胞腔链复形。}

与同伦群的情形一样，对复形对 $(K,L)$ 可以自然地定义相对同调群与相对上同调群，其中 $K$ 是单纯复形，$L$ 是它的子复形。

对 $(K,L)$ 的\emph{相对链}，就是 $K$ 的链按「两两之差为 $L$ 的链」这一等价关系所得的等价类；也就是说，两条链若模 $L$ 的链互相同余，便视为同一条。于是得到\emph{相对链复形}
\begin{equation}
  C_j(K,L)=C_j(K)/C_j(L). \tag{3.1}
\end{equation}
$C(K,L)$ 上的边缘算子 $\partial$，由 $K$-链上通常的边缘算子诱导得到；由于 $\partial C_j(L)\subset C_{j-1}(L)$，这确实给出一个良定义的算子：
\[
  \cdots\to C_j(K,L)\xrightarrow{\partial}C_{j-1}(K,L)\xrightarrow{\partial}\cdots .
\]
对 $(K,L)$ 的第 $j$ 个整相对同调群定义为
\begin{equation}
  H_j(K,L;\Z)=\Ker\partial/\Image\partial. \tag{3.2}
\end{equation}
这里 $\Ker\partial$ 与 $\Image\partial$ 的元素，分别称为\emph{相对闭链}与\emph{相对边缘}。

系数取任意交换群 $G$ 的相对同调群 $H_j(K,L;G)$，仿照以 $G$ 为系数的绝对同调群的做法，用复形 $C(K,L)\otimes G$ 定义；而 $G$ 上的相对上同调群 $H^j(K,L;G)$，则用复形 $\Hom(C(K,L),G)$ 及相对上边缘算子 $\partial^*$ 定义。相对 $G$-上链在几何上可以看作一个取值为 $G$ 的函数，它定义在 $K$ 的单形集合上，并在 $L$ 的单形上取零。

与同伦群的情形十分相似，可以得到对 $(K,L)$ 的同调与上同调的长正合列：
\begin{equation}
  \cdots\xrightarrow{i_*}H_j(K)\xrightarrow{j_*}H_j(K,L)\xrightarrow{\partial}
  H_{j-1}(L)\xrightarrow{i_*}H_{j-1}(K)\to\cdots, \tag{3.3}
\end{equation}
\begin{equation}
  \cdots\xleftarrow{i_*}H^j(K)\xleftarrow{j_*}H^j(K,L)\xleftarrow{\partial^*}
  H^{j-1}(L)\xleftarrow{i_*}H^{j-1}(K)\leftarrow\cdots . \tag{3.4}
\end{equation}

% [原文 p.58]
相对同调群与相对上同调群的构造，以及相应的正合列 (3.3)、(3.4)，对复形对范畴而言都是函子性的；该范畴的态射是复形对的单纯映射 $f:(K_1,L_1)\to(K_2,L_2)$。这些对象全都是同伦不变量；诱导同态 $f_*$、$f^*$ 不受复形对映射的同伦影响。

用类似的作法还可以得到更一般的三重 $(K,L,M)$（$K\supset L\supset M$）的正合列：
\begin{equation}
  \cdots\xrightarrow{i_*}H_j(K,M)\xrightarrow{j_*}H_j(K,L)\xrightarrow{\partial}
  H_{j-1}(L,M)\xrightarrow{i_*}H_{j-1}(K,M)\to\cdots, \tag{3.5}
\end{equation}
\begin{equation}
  \cdots\xleftarrow{i_*}H^j(K,M)\xleftarrow{j_*}H^j(K,L)\xleftarrow{\partial^*}
  H^{j-1}(L,M)\xleftarrow{i_*}H^{j-1}(K,M)\leftarrow\cdots . \tag{3.6}
\end{equation}

注意到复形 $K$ 对子复形 $L$ 的商空间 $\bar K=K/L$ 具有自然的三角剖分，便得到下述结果（\emph{切除定理}）：
\begin{equation}
  \begin{aligned}
    H_j(K,L)&\cong H_j(K/L),\qquad j>0,\\
    H_0(K,L)&\cong 0\qquad(\text{$K$ 连通且 }L\neq\varnothing).
  \end{aligned} \tag{3.7}
\end{equation}
% [Tutor Note: (3.7) 的 $j>0$ 一栏里 $L$ 不必连通，也不必非空；而 $H_0$ 一栏的 $L\neq\varnothing$ 不可省——$L=\varnothing$ 时 $H_0(K,\varnothing)=H_0(K)$ 当然不消失。记 $K/L$ 时把 $L$ 压成一点，正是这个「压点」给出切除。]
切除定理来自这样一个观察：商空间 $K/L$ 显然同伦等价于把 $L$ 上的锥沿 $L$ 粘到 $K$ 上所得的复形：
\[
  K/L\simeq K\cup_L CL.
\]

对悬垂 $\Sigma K=(CK)_1\cup_K(CK)_2$，由对 $(\Sigma K,(CK)_2)$ 的同调正合列，再结合复形 $CK$ 的可缩性，可以推出下述重要的\emph{悬垂同构}：
\begin{equation}
  H_j(K,x_0)\cong H_{j+1}(\Sigma K,x_0),\qquad j\ge 0. \tag{3.8}
\end{equation}
从离散两点空间 $K=S^0=\{x_0,x_1\}$ 的同调出发，反复使用悬垂同构，用归纳法便得到球面的同调群：
\begin{equation}
  \begin{aligned}
    H_j(S^n,x_0)&=0,\quad j\neq n,\\
    H_n(S^n,x_0;\Z)&\cong H^n(S^n,x_0;\Z)\cong\Z,\\
    H_n(S^n,x_0;G)&\cong H^n(S^n,x_0;G)\cong G.
  \end{aligned} \tag{3.9}
\end{equation}

按 Eilenberg--Steenrod 的一个定理，同调（与上同调）论由下列条件（\emph{Eilenberg--Steenrod 公理}）唯一确定（即在同构意义下唯一）：函子性、同伦不变性、对的正合列 (3.3) 的存在性，以及正规化公理（本质上即性质 (3.9)）；此外还要求切除公理 (3.7) 成立。
% [原文 p.59]
正是后一条公理把同调论与同伦论区分开来。

20 世纪 60 年代起，所谓\emph{广义}（或\emph{超常}）同调（与上同调）论开始被大量用作拓扑学的工具。这些理论与通常的同调论只差一条正规化公理：球面 $(S^n,s_0)$ 的\emph{广义}同调群可以在 $n$ 以外的维数上非平凡；等价地说，可缩空间（例如圆盘或单点）的（绝对）广义同调群可以在非零维数（包括负维数）上非平凡。下面我们将讨论与此有关的问题。

Poincaré 对偶可以推广到相对的情形，而不必改变底层的几何假设。设 $M^m$ 是闭的、可定向的 PL 流形，$K\subset M^m$ 是子复形。开补集 $U=M^m\setminus K$ 总有有限子复形 $L\subset U$ 作为它的形变收缩核。我们把对 $(M^m,U)$ 的同调理解为对 $(M^m,L)$ 的同调。于是得到 Poincaré 对偶的下述推广形式：
\begin{equation}
  H_j(K)\cong H^{m-j}(M^m,U). \tag{3.10}
\end{equation}

\begin{example}{1}
取 $M^m=S^m$，即 $m$ 维球面。只要 $S^m\setminus U$ 非空，三元组 $(S^m,U,*)$ 的上同调正合列（见式 (3.5)）就给出同构
\[
  H^{m-j}(S^m,U)\cong H^{m-j-1}(U,*),\qquad j\neq 0.
\]
% [Tutor Note: 原书 p.59 此处引「式 (3.5)」，但这里讲的是上同调三元组的正合列，对应的是式 (3.6)（(3.5) 是同调版）；疑为交叉引用笔误，正文照抄。]
由此可以推出 Alexander 对偶的同构——见式 (2.24)（$j=0$ 时也成立）：
\begin{equation}
  H_j(K,x_0)\cong H^{m-j-1}(U,u_0),\qquad x_0\in K,\quad u_0\in U. \tag{3.11}
\end{equation}
\end{example}

\begin{example}{2}
设 $M^m$ 是任意可定向流形，$K=W^m$ 是 $M^m$ 中某个非空开集的闭包。于是 $K$ 是 $M^m$ 中同维数 $m$ 的子流形，带边界 $\partial W^m$。由式 (3.11) 与切除定理 (3.7) 便直接得到同构
\[
  \begin{aligned}
    H_j(W^m,\partial W^m)&\cong H^{m-j}(W^m),\\
    H^j(W^m,\partial W^m)&\cong H_{m-j}(W^m).
  \end{aligned}
\]
这些同构就是所谓 \emph{Lefschetz 对偶}。
\end{example}

同伦不变的同调论，最方便是在可数胞腔复形（或 CW 复形）的范畴中构造和研究。我们已经遇到过——例如在构造单纯复形的 Poincaré 对偶的过程中——这种必要性，或者至少是方便之处——考虑把单纯复形合并，使结构单元不再是单形，而是例如凸多面体或多胞形。
% [原文 p.60]
CW 复形是这类对象中最一般、也最方便的一种。

先设 $K$ 是一个单纯复形。若给定一族各种维数 $n$ 的胞腔（圆盘）$D_\alpha^n$，以及一组单纯映射
\[
  \psi_\alpha^n: D_\alpha^n\to K,
\]
并满足下述性质，我们就说 $K$ 被\emph{分解}（或\emph{剖分}）成胞腔：
\begin{itemize}
  \item[(i)] 每个映射 $\psi_\alpha^n$ 在胞腔 $D_\alpha^n$ 的内部 $\operatorname{Int}D_\alpha^n$ 上是一一的；
  \item[(ii)] 各胞腔（即各像 $\psi_\alpha^n(D_\alpha^n)$）之并是整个复形 $K$，且它们的内部 $\psi_\alpha^n(\operatorname{Int}D_\alpha^n)$ 两两不交；
  \item[(iii)] 对每个 $n$，像 $\psi_\alpha^n(\partial D_\alpha^n)$（即胞腔 $\psi_\alpha^n(D_\alpha^n)$ 的边界）落在维数小于 $n$ 的胞腔之并中。
\end{itemize}

在这些条件下，我们说 $K$ 被赋予了 CW 复形 $\hat K$ 的结构。由这个定义立刻看出：给定 $K$ 的一个定向（假定 $K$ 可定向），这种胞腔剖分的每个 $n$ 维胞腔 $\psi_\alpha^n(D_\alpha^n)$，都是该胞腔内所含 $K$ 的（定向）$n$-单形的线性组合：
\[
  \psi_\alpha^n(D_\alpha^n)=\sum_\gamma\lambda_{\alpha\gamma}\sigma_\gamma^n,
  \qquad \lambda_{\alpha\gamma}=0,\pm1,
\]
其中 $\sigma_\gamma^n$ 是 $K$ 的 $n$ 维单形。（注意：不同胞腔即使相交也只沿边界相交，因此每个 $n$-单形 $\sigma_\gamma^n$ 的内部恰好含于一个胞腔的内部。）把胞腔 $\psi_\alpha^n(D_\alpha^n)$ 记作 $\kappa_\alpha^n$。每个圆盘 $D_\alpha^n$ 的边界 $\partial D_\alpha^n$ 被单纯地映入 $K$ 中维数小于 $n$ 的单形之并，即 $K$ 的 $(n-1)$-骨架；因此，§2 中已对单纯链定义的边缘算子 $\partial$，就给出 $\hat K$ 的每个 $n$-胞腔的边缘 $\partial\kappa_\alpha^n$，作为 $n-1$ 维胞腔的整系数线性组合的表达式：
\begin{equation}
  \partial\kappa_\alpha^n=\sum_\delta[\kappa_\alpha^n:\kappa_\delta^{n-1}]\,\kappa_\delta^{n-1},
  \qquad \partial\circ\partial=0. \tag{3.12}
\end{equation}

整数 $[\kappa_\alpha^n:\kappa_\delta^{n-1}]$ 称为 CW 复形 $\hat K$ 的\emph{关联数}。

CW 复形 $\hat K$ 的\emph{胞腔同调群}与\emph{胞腔上同调群}，如今仿照单纯的情形来定义：用整胞腔链（它们构成 $C_*(\hat K)$）与胞腔上链；系数取任意交换群 $G$ 时，则改用 $C_*(\hat K)\otimes G$ 与 $\Hom(C_*(\hat K),G)$。一个 CW 复形的胞腔同调群与上同调群，与其单纯对应物（即上面 §2 中所定义者）一致，也就是说彼此自然同构。

% [原文 p.61]
CW 复形 $K$ 的 $n$-\emph{骨架} $K^n$，是维数 $\le n$ 的胞腔之并。显然每个商空间 $K^n/K^{n-1}$ 都是球面的\emph{花束}：
\[
  K^n/K^{n-1}\cong S^n_1\vee\cdots\vee S^n_m\vee\cdots,
\]
其中的球面依次代表各个 $n$-胞腔 $\kappa_\alpha^n$，$\alpha=1,2,\dots,m,\dots$。于是群 $H_n(K^n,K^{n-1};\Z)$ 是自由交换群，其自由生成元由各胞腔 $\kappa_\alpha^n$ 代表；边缘算子 $\partial$ 也可以另作定义：它是对 $(K^n,K^{n-1})$ 与 $(K^{n-1},K^{n-2})$ 这两对的正合同调列中所出现的映射 $\partial_*$、$j_*$ 的复合：
\[
  H_n(K^n,K^{n-1};\Z)\xrightarrow{\partial_*}H_{n-1}(K^{n-1};\Z)
  \xrightarrow{j_*}H_{n-1}(K^{n-1},K^{n-2};\Z).
\]

CW 复形的定义可以稍加放宽，不再依赖一开始给定的单纯剖分。所谓拓扑空间 $X$ 是\emph{CW 复形}，是指给定一族各种维数的圆盘（球体）$D_\alpha^n$ 上的映射 $\psi_\alpha^n:D_\alpha^n\to X$（它们的像，严格说是这些映射本身，称为\emph{胞腔}），并满足与前面相同的性质 (i)、(ii)、(iii)：每个 $\psi_\alpha^n$ 在 $\operatorname{Int}D_\alpha^n$ 上是一一的；各胞腔之并是整个 $X$；各胞腔的内部 $\psi_\alpha^n(\operatorname{Int}D_\alpha^n)$ 两两不交；对每个 $n$，集合 $\psi_\alpha^n(\partial D_\alpha^n)=\psi_\alpha^n(S^{n-1})$ 含于维数 $\le n-1$ 的胞腔之并中。此外还要求 $X$ 上的拓扑是相应的\emph{粘合拓扑}：映射 $f:X\to Z$（$Z$ 为任意拓扑空间）连续，只要各复合映射 $f\circ\psi_\alpha^n$ 都连续。换言之，只要 $f$ 是各胞腔上映射的并，且这些映射两两在胞腔边界的公共部分上一致，$f$ 就连续。

在这一般情形下，CW 复形 $X$ 的 $n$-骨架是由维数 $\le n$ 的胞腔构成的子复形 $X^n$。显然每个商空间 $X^j/X^{j-1}$ 仍是 $j$ 维球面的花束：
\[
  X^j/X^{j-1}\cong S^j_1\vee\cdots\vee S^j_m\vee\cdots,
\]
其中每个球面 $S^j_\alpha$ 恰好来自一个 $j$-胞腔 $\kappa_\alpha^j$（$=\psi_\alpha^j(D_\alpha^j)$），$\alpha=1,\dots,m,\dots$。边缘算子 $\partial$ 具有与前面相同的形式（见式 (3.12)）；只是现在，一对（定向）胞腔 $\kappa_\alpha^{j+1}$、$\kappa_\beta^j$ 的\emph{关联数} $[\kappa_\alpha^{j+1}:\kappa_\beta^j]$，改定义为复合映射
\[
  \partial(D_\alpha^{j+1})\to\partial(\kappa_\alpha^{j+1})\to S_\beta^j
\]
的度，该复合映射由 $\psi_\alpha^{j+1}$ 与上述粘合确定。这本质上就是复合映射 $\psi_\alpha^{j+1}|_{\partial(D_\alpha^{j+1})}$ 与粘合映射 $X^j/X^{j-1}\to\vee S_\alpha^j$ 之复合，把 $j$ 维球面 $\partial(D_\alpha^{j+1})$ 绕到花束中相应的球面 $S_\beta^j$ 上，所绕的（带符号的）次数。由此推出 $\partial\circ\partial=0$；随后按通常方式定义 CW 复形 $X$ 的同调群与上同调群，得到的群与前面定义的群自然同构。

% [原文 p.62]
任何 CW 复形都可由反复施行下述操作来构造。给定空间 $X$ 与一列可数个映射 $f_\alpha:S^{j-1}\to X$，把各个圆盘 $D_\alpha^j$ 的边界经 $f_\alpha$ 与相应的 $f_\alpha(S^{j-1})$ 等同，即把圆盘 $D_\alpha^j$ 粘到 $X$ 上，所得的新空间记为 $Y$：
\begin{equation}
  Y=X\cup_{(f_1,\dots,f_\alpha,\dots)}\bigl(D_1^j\cup\cdots D_\alpha^j\cup\cdots\bigr). \tag{3.13}
\end{equation}
$Y$ 上的拓扑按自然方式定义。于是：从可数离散的点集 $K^0$（它最终成为 $0$-骨架）出发，先粘上 $1$-胞腔（$j=1$），再向所得的 $1$-骨架粘上 $2$-胞腔，如此继续。

最后我们指出一个事实：每个 CW 复形都同伦等价于某个单纯复形。

\begin{example}{1}
球面 $S^n$ 最简的胞腔剖分如下：在唯一一个 $0$ 维胞腔（顶点）$\kappa^0$ 上，用显然的映射 $\psi:D^n\to S^n$ 粘上一个 $n$ 维胞腔 $\kappa^n$，该映射把 $D^n$ 的边界映到点 $\kappa^0$：$\psi(\partial D^n)=\{\kappa^0\}$。显然 $\partial\kappa^0=0=\partial\kappa^n$，于是同调群为：
\[
  \begin{gathered}
    H_0(S^n)\cong\Z\cong H_n(S^n),\\
    H_j(S^n)=0\quad\text{当 }j\neq0,n.
  \end{gathered}
\]
\end{example}

\begin{example}{2}
a) 实射影空间有如下胞腔剖分，边缘算子如式中所标：
\begin{align}
  \RP^n&=\kappa^0\cup\kappa^1\cup\cdots\cup\kappa^n,\notag\\
  \kappa^j&=\RP^j\setminus\RP^{j-1}
    =\{(x^0:\dots:x^n)\mid x^j\neq0,\ x^k=0\ (k>j)\},\notag\\
  \partial\kappa^0&=\partial\kappa^1=\partial\kappa^3=\partial\kappa^5=\cdots
    =\partial\kappa^{2k+1}=\cdots=0,\quad
    \partial\kappa^{2k}=2\kappa^{2k-1},\ k>0. \tag{3.14}
\end{align}
b) 复射影空间的情形为：
\begin{align}
  \CP^n&=\kappa^0\cup\kappa^2\cup\cdots\cup\kappa^{2n},\notag\\
  \kappa^{2j}&=\CP^j\setminus\CP^{j-1}
    =\{(z^0:\dots:z^n)\mid z^j\neq0,\ z^k=0\ (k>j)\},\notag\\
  \partial\kappa^{2j}&=0,\quad j>0. \tag{3.15}
\end{align}
\end{example}

\begin{example}{3}
闭的可定向曲面与不可定向曲面，是由把凸多边形的适当边对等同而构造的（见图 2.22、2.23）；由此得到这些曲面的如下极小胞腔剖分：

% [原文 p.63]
(i) 对可定向曲面 $M_g^2$：
\[
  \kappa^0\cup\kappa_1^1\cup\cdots\cup\kappa_{2g}^1\cup\kappa^2,\qquad
  \partial\kappa_j^1=0,\quad j=1,\dots,2g,\qquad \partial\kappa^2=0.
\]
(ii) 对第一族不可定向曲面 $N_{1,g}^2$，$g\ge1$：
\[
  \kappa^0\cup\kappa_1^1\cup\cdots\cup\kappa_{2g}^1\cup\kappa^2,\qquad
  \partial\kappa_j^1=0,\quad j=1,\dots,2g,\qquad \partial\kappa^2=2\kappa_0^1.
\]
% [Tutor Note: 原书 p.63 在 (ii) 中列出的 1-胞腔是 $\kappa_1^1,\dots,\kappa_{2g}^1$，却写 $\partial\kappa^2=2\kappa_0^1$，而 $\kappa_0^1$ 并未出现在这族胞腔里；对照 (iii) 可知此处疑为笔误，应为 $2\kappa_1^1$。正文照抄。]
(iii) 对第二族不可定向曲面 $N_{2,g}^2$，$g\ge0$：
\[
  \kappa^0\cup\kappa_0^1\cup\cdots\cup\kappa_{2g}^1\cup\kappa^2,\qquad
  \partial\kappa_j^1=0,\quad j=0,\dots,2g,\qquad \partial\kappa^2=2\kappa_{2g}^1.
\]

由这些公式容易算出闭曲面的同调群。

这些曲面上 $1$ 维闭链的相交指数（对不可定向曲面按模 $2$ 取）由上述基本 $1$-闭链两两的相交指数确定，具体如下：记 $a_i=\kappa_{2i}^1$，$b_i=\kappa_{2i-1}^1$，则对所有曲面 $M_g^2$、$N_{1,g}^2$、$N_{2,g}^2$ 都有
\begin{equation}
  a_i\circ a_j=b_i\circ b_j=0,\qquad a_i\circ b_j=\delta_{ij}, \tag{3.16}
\end{equation}
其中对不可定向曲面，$\delta_{ij}$ 按模 $2$ 理解。对曲面 $N_{2,g}^2$ 还存在附加的基本 $1$-闭链 $\kappa_0^1=c$；由这条 $1$-闭链产生的附加基本相交指数如下：
\begin{equation}
  c\circ c=1,\qquad c\circ a_i=c\circ b_i=0\quad(\text{mod }2). \tag{3.17}
\end{equation}
满足 (3.16)（以及适当情形下的 (3.17)）的任意一族闭链 $a_i$、$b_i$（$i=1,\dots,g$；对 $N_{2,g}^2$ 还要加上 $c$），称为该曲面上闭链的\emph{典范基}。

借助 Poincaré 对偶的同构——它在每种情形下都把群 $H_0$ 的基本元映到 $H^2$ 的基本上同调类——便可由 (3.16)（对 $N_{2,g}^2$ 还要用 (3.17)）推出上同调环
\[
  H^*(M_g^2;\Z),\qquad H^*(N_{1,g}^2;\Z/2),\qquad H^*(N_{2,g}^2;\Z/2)
\]
的结构。
\end{example}

前面已经提到过 \emph{Hopf 定理}：
\begin{equation}
  H_1(K;\Z)\cong\pi_1(K)/[\pi_1(K),\pi_1(K)], \tag{3.18}
\end{equation}
它对任何 CW 复形 $K$ 都成立。

% CHUNK ch03c DONE
